\documentclass[11pt]{article}
\usepackage[a4paper,margin=25mm]{geometry}
\usepackage{amsmath,amssymb,amsthm}
\usepackage[hidelinks]{hyperref}
\emergencystretch=2em
\newtheorem{theorem}{Theorem}
\title{A regularity exponent cannot have slope tending to\newline
minus infinity at infinity\\\large Disproof of Conjecture 00000007991}
\author{gaochengzhecpu}\date{}
\begin{document}\maketitle
\begin{abstract}
Every real function bounded below on a tail is incompatible with an
ordinary derivative tending to $-\infty$ as its argument tends to
$+\infty$. A mean value theorem argument proves this directly.
An optimal H\"older regularity exponent is nonnegative, so the
claimed slope limit for $\alpha(p)$ at $p=\infty$ is impossible.
The argument applies to every such exponent function and requires
no formula for the actual optimal $p$-Laplace exponent.
\end{abstract}

\section*{A general calculus obstruction}
\begin{theorem}
Let $f$ be a real function bounded below on a tail of the real line.
Its ordinary derivative cannot satisfy
\[
             \lim_{x\to+\infty} f'(x)=-\infty.
\]
In particular, no eventually nonnegative real function has this
ordinary derivative limit.
\end{theorem}
\begin{proof}
Suppose such a limit exists. There are real numbers $L$ and $a$ such
that, for every $x\geq a$, both $f(x)\geq L$ and $f'(x)\leq-1$.
Here and below an ordinary derivative on this tail is meant.
Set
\[
                 b=a+f(a)-L+1.
\]
Since $f(a)\geq L$, we have $b>a$. Differentiability on the tail
implies continuity on $[a,b]$ and differentiability on $(a,b)$.
The real mean value theorem gives a point $c\in(a,b)$ with
\[
 \frac{f(b)-f(a)}{b-a}=f'(c)\leq-1.
\]
Consequently
\[
 f(b)\leq f(a)-(b-a)=L-1<L,
\]
contradicting the lower bound at $b$.
\end{proof}

The proof only needs the eventual slope bound $f'\leq-1$.
It concerns derivatives on actual real intervals, rather than a
discrete sequence or finite-difference surrogate.

\newpage
\section*{Application to the exact conjecture}
The source defines $\alpha(p)$ as the optimal exponent of
$C^{1,\alpha}$ regularity of solutions of
\[
            \operatorname{div}(|\nabla u|^{p-2}\nabla u)=0.
\]
One of its explicit assertions, in both languages, is that the
slope of the graph of $\alpha(p)$ tends to $-\infty$ as
$p\to\infty$.

A H\"older regularity exponent is nonnegative by its definition.
For the usual range $0<\alpha\leq1$, a local gradient estimate has
the form
\[
 |\nabla u(x)-\nabla u(y)|\leq C|x-y|^\alpha.
\]
An optimal exponent, whether attained or defined as a supremum of
admissible nonnegative exponents, therefore cannot be negative.
Only this necessary range condition is used here. In particular,
the exact value, existence of an optimal endpoint estimate, and
dependence on the spatial dimension are immaterial to the argument.

Apply Theorem 1 with $f(p)=\alpha(p)$ and $L=0$ on the tail $p>2$.
It gives
\[
 \neg\left(\alpha(p)\geq0\text{ for all }p>2
       \quad\text{and}\quad
       \alpha'(p)\longrightarrow-\infty
                     \text{ as }p\to+\infty\right).
\]
Thus there is no candidate exponent function meeting even these
two necessary parts of the conjecture. Additional PDE conditions
can only restrict that already empty class of candidates.

\section*{Scope and derivative convention}
This refutes the negative-infinite-slope conjunct, hence the
conjunction in the source. It does not purport to compute the
actual optimal exponent, validate or refute the proposed
eigenvalue formula separately, or settle the behavior as
$p\to1$. No interpretation or sign assumption concerning the
unspecified eigenvalue $\lambda_0$ is needed.

The asserted graph slope is interpreted as the ordinary real
derivative with respect to the stated horizontal coordinate $p$.
It is not a slope after reparametrizing the axis, nor a limit taken
only along a selected discrete sequence of values of $p$.
The theorem is about a whole tail of real arguments.

In Lean, \texttt{deriv} is a total function and equals zero at a
point where the function is not differentiable. A proposed limit
to $-\infty$ makes its values eventually at most $-1$. This
\emph{forces} differentiability at every point of that tail;
the proof explicitly derives this fact and then uses the
ordinary mean value theorem. No prior differentiability
hypothesis is silently added to the final Lean theorem.

\newpage
\section*{Formal conclusions and reproduction}
The Lean project also checks the source's displayed scalar formula
when its eigenvalue symbol is a fixed real parameter $\ell$.
Nonnegativity at $p=3$ forces $\ell>-2$: otherwise its first
branch $2/(4+3\ell)$ is negative.
Put $c=\ell+2>0$ and
$A=\max(3,\ell+5,4/c)$. For $p>A$, both denominators
are positive and
\[
 (p-1)^2-(cp-2)=p(p-\ell-4)+3>0.
\]
Thus the minimum in the formula is exactly $1/(p-1)$
on that entire open tail. Its actual derivative is
$-1/(p-1)^2\to0$, also incompatible with the claimed
limit. This supplementary result proves eventual
differentiability for the displayed constant-parameter
formula itself. The main contradiction on the preceding
pages does not require that formula or any assumption
about whether an intended eigenvalue could depend on $p$.

The final theorem \texttt{conjecture\_cusp\_impossible}
directly excludes every graph satisfying the necessary
nonnegative range and the claimed ordinary slope limit.
The general results use the actual real mean value theorem
on the interval constructed above. A separate version takes
values directly in the nonnegative reals. No PDE or special
formula is built into the general function hypotheses.

Lean 4.19.0 and Mathlib revision\\
\texttt{c44e0c8ee63ca166450922a373c7409c5d26b00b}
are pinned. From the submitted \texttt{lean/} directory:
\begin{verbatim}
lake exe cache get
lake build
lake env lean -DwarningAsError=true Main.lean
\end{verbatim}
The complete project is checked by a fresh build and strict
Lean checking, with warnings treated as errors. Printed
theorem dependencies are audited for standard foundational
axioms only. There are no proof placeholders, custom axioms,
or numerical computation oracles. No auxiliary numerical
program is needed; the argument is an exact calculus proof.

Compile this document with \texttt{tectonic main.tex}.
The submission includes the exact bilingual source,
provenance, actual build logs, artifact hashes and PDF
checks. An author agent developed the formal proof and
supplied a self-review; the parent reviewed the proof and
the final artifacts separately. No external independent
review or novelty of the elementary calculus fact is claimed.

\begin{thebibliography}{9}
\bibitem{source} The Last Math Competition,
Conjecture 00000007991. Exact bilingual source and
provenance included in this submission.
\bibitem{mathlib} The Mathlib community,
\emph{Mathlib}, revision specified above,
the calculus library for real mean value theorems.
The submitted source explicitly applies
\texttt{exists\_deriv\_eq\_slope} on a real interval.
\end{thebibliography}
\end{document}
